\section{Complete CALM and Complete CAP}
\label{sec:complete-calm}

The previous two sections developed orthogonal extensions of Functional CALM:
Lamport histories give Consistency a distributed reading, while relational
specifications add output choice.  We now combine them.  First we identify
coordination with commitment and state when conflicting commitment sequences
cannot coexist; then Complete CALM equates Consistency, Lifted Monotonicity, and
Distributed Consistency, and Complete CAP states the resulting partition
boundary.

Figure~\ref{fig:complete-formalism-summary} recaps the distributed and
relational extensions combined here.

\begin{figure*}[t]
  \centering
  \small
  \setlength{\tabcolsep}{6pt}
  \renewcommand{\arraystretch}{1.25}
  \begin{tabular}{@{}l l l@{}}
    \toprule
    \textbf{Distributed histories} & & \\
    \midrule
    participant; message & participant holds one execution state;
      $m=(i_m,s_m)$ carries input and strategy components & Sec.~\ref{sec:histories} \\
    history; consistent cut & locally ordered events plus matched sends and receives;
      cut closed under happens-before & Def.~\ref{def:distributed-history} \\
    projection; endpoint & $\pi_p(C)$: transitions affecting $p$;
      $\operatorname{end}(H,C)=(\gamma_p(C))_{p\in P}$ & Def.~\ref{def:semantic-projection} \\
    history confluence & at every common input $j$, all participants can reach one stable execution state $((j,R),\mathsf{none})$ & Sec.~\ref{sec:projection-confluence} \\
    \addlinespace
    \textbf{Output choice} & & \\
    \midrule
    fiber at $i$ & $S(i)=\{o:(i,o)\in S\}$ & Eq.~\ref{eq:fiber} \\
    Lifted Monotonicity & $i\InputOrd j\Rightarrow S(i)\JointOutputOrd S(j)$; finite sets in $S(i)$ have common refinements in $S(j)$ & Def.~\ref{def:lifted-monotonicity} \\
    output retirement & removed $(k,x)$ is covered by retained $(k,u)$ with $x\OutputOrd u$ & Def.~\ref{def:output-retirement} \\
    exposable output & $\mathsf{prepare}(o)$ is the identity; exposure retires only dominated rows & Def.~\ref{def:exposable-relational} \\
    output-prepare commitment & removes undominated conflicts while leaving their input valid & Sec.~\ref{sec:safe-exposure} \\
    \addlinespace
    \textbf{Results} & & \\
    \midrule
    Complete CALM & Consistency $\iff$ Lifted Monotonicity $\iff$ Distributed Consistency & Thm.~\ref{thm:complete-calm} \\
    coordination & application of a commitment & Def.~\ref{def:coordination} \\
    conflicting sequences & cannot coexist without mutual exclusion & Thm.~\ref{thm:commitments-require-coordination} \\
    \bottomrule
  \end{tabular}
  \caption{The distributed and relational extensions combined by Complete
  CALM, for reference.}
  \label{fig:complete-formalism-summary}
\end{figure*}

\subsection{Coordination}
\label{sec:coordination}

Consistency is confluence: independently evolving states can still meet.
Coordination is its operational counterpart: the system makes a binding choice
that removes some otherwise possible behavior.  In our transition system, that
act is exactly a commitment.  A mutex, lock, atomic ordering mechanism, or
agreement protocol is in scope when it serves this role.  Computation or
communication that removes no such alternative is not coordination merely
because it uses one of these mechanisms.

\begin{definition}[Coordination]
\label{def:coordination}
A system execution \emph{coordinates} when a preparation expansion applies a
commitment: it irrevocably removes a row that neither input nor output retirement
accounts for.  It is \emph{coordination-free} when every preparation is the
identity; exposure events remain explicit and may perform output retirements.
\end{definition}

Hereafter, \emph{reachable} means reachable without commitment.  A system
selects an output through its exposure sequence.  When the output is exposable,
preparation is the identity; otherwise its preparation expansion includes a
commitment before exposing the prepared output.

In a shared-memory system, a mutex or atomic instruction can establish a
commitment by fixing an order or choice.  In a distributed system, a commitment
that conflicts with an alternative that must remain independently selectable
requires mutual exclusion before it is applied; agreement is one way to
establish that exclusion.

\begin{theorem}[Conflicting commitment sequences cannot coexist]
\label{thm:commitments-require-coordination}
Suppose one participant can apply a commitment $c$ at a reachable state
$\sigma=(i,R)$, producing $(i,R')=c(\sigma)$ and removing some $x\in R(i)$ by
commitment: no output in $R'(i)$ refines $x$.  Suppose a second participant can
independently select $x$ after receiving the same input events, obtaining a
state $t$ whose fiber at $i$ contains only refinements of $x$.  Then $c(\sigma)$
and $t$ cannot be brought to a common state with current input $i$.
\end{theorem}

\begin{proof}
By hypothesis, the second participant can reach $t$, whose fiber at $i$ is
contained in $\{u\in\fiber{i}:x\OutputOrd u\}$.  The first participant's fiber
$R'(i)$ contains no such $u$.  Their fibers at $i$ are therefore disjoint, so
the two states cannot be brought to a common state whose current input is $i$.
Neither preparation nor an input event can restore a removed row or return to
$i$.
\end{proof}

If both sequences must remain independently available, a system must mutually
exclude them before either participant commits; distributed agreement is one
way to do so.  Static authority or a predetermined strategy can instead make
one sequence unavailable, in which case no runtime agreement follows from the
theorem.

\subsection{The Complete CALM theorem}
\label{sec:complete-calm-theorem}

Complete CALM turns the coordination boundary into a property of the
specification.  It equates Consistency---confluence of runs under the canonical
commitment basis---with Lifted Monotonicity.  The distributed execution model
supplies the equivalent third view, Distributed Consistency.

\begin{theorem}[Complete CALM]
\label{thm:complete-calm}
For an ordered relational specification, the following are equivalent:
\begin{enumerate}
  \item Consistency;
  \item Lifted Monotonicity;
  \item Distributed Consistency.
\end{enumerate}
\end{theorem}

\begin{proof}
Items~(1) and~(3) are equivalent by
Proposition~\ref{prop:distributed-consistency}.  It remains to prove the
substantive equivalence of items~(1) and~(2).

Proposition~\ref{prop:confluence} says that Consistency means: for every
finite set of runs ending at $(i_r,R_r)$ and every common refinement
$j$ of the $i_r$, some output $u\in\fiber{j}$ survives in every run,
$u\in\bigcap_rR_r(j)$.  By Lemma~\ref{lem:run-residual}, this holds exactly
when $u$ refines every output exposed on those runs; each such output $o$ was
exposed at some $i_o\InputOrd j$ with $o\in\fiber{i_o}$.

\emph{(2) $\Rightarrow$ (1): Lifted Monotonicity implies Consistency.}
Assume Lifted Monotonicity.  The finitely many rows $(i_o,o)$ exposed on the
runs have common refinement $j$, so
Proposition~\ref{prop:lifted-monotonicity-decomposition}(3) supplies
$u\in\fiber{j}$ refining them all.  Thus $u$ survives every run.

\emph{(1) $\Rightarrow$ (2): Consistency implies Lifted Monotonicity.}
Let rows $(i_1,o_1),\ldots,(i_k,o_k)$ have a common refinement $j$.  For each
$r$, start a run at $\gamma_0(i_r)$ and complete
$\mathsf{prepare}(o_r);\mathsf{expose}(o_r)$ under the canonical commitment
basis; then
$R_r(j)=\{u\in\fiber{j}:o_r\OutputOrd u\}$.  Consistency supplies an output
$u$ that survives every run, so it refines every $o_r$.  This is form~(3) of
Proposition~\ref{prop:lifted-monotonicity-decomposition}, hence Lifted
Monotonicity.

\end{proof}



The theorem gives a specification-level test for history confluence and closes
the operational loop.  Exposability follows rather than being assumed.  In a
consistent specification, every $o\in R(i)$ is exposable at every reachable
$(i,R)$: for valid $j\succeq_I i$ and $x\in R(j)$, confluence of the run
extended by the exposure sequence for $o$ with the run extended by
$\mathsf{input}(j)$ and the exposure sequence for $x$ supplies $u\in R(j)$
above both outputs; then
Proposition~\ref{prop:exposability} applies.  Thus Lifted Monotonicity makes all
reachable outputs exposable, so progress needs no commitment and is
coordination-free.  Its failure exhibits a finite family without a common
refinement, so selecting an offending output requires a commitment and hence
coordination.  When different participants must retain independently
selectable conflicting alternatives,
Theorem~\ref{thm:commitments-require-coordination} shows that their sequences
must be mutually excluded; distributed agreement is one mechanism for doing so.

\subsection{The Complete CAP Theorem}
\label{sec:cap-corollary}

The coordination boundary becomes an availability boundary during a
partition: participants in different components cannot establish a new shared
commitment.  A \emph{partitioned history} divides participants into two or more
components and has no cross-component message receive while the partition
lasts, though participants may still receive environment input and expose
outputs.  At a cut, first complete any pending exposure; this yields one stable
execution state per participant.  We use an interface-wide form of
\emph{Availability}: every output in a participant's current fiber at that
stable state can complete its exposure sequence with identity
preparation---that is, without commitment or a cross-component wait.  This is
stronger than requiring some response to one operation, because the full
declared interface treats every permitted output as available.
Appendix~\ref{app:deployment} restricts the condition to outputs a
participant is authorized to expose.  The history has \emph{Partition
Consistency} when its commitment-free participant
projections are confluent at every common refinement of their current inputs.

\begin{theorem}[Complete CAP]
\label{thm:complete-cap}
For an ordered relational specification $\Spec$, the following are equivalent:
\begin{enumerate}
  \item every consistent cut of every commitment-free partitioned
        history has Availability and Partition Consistency;
  \item $\Spec$ satisfies Lifted Monotonicity.
\end{enumerate}
\end{theorem}

\begin{proof}
Suppose first that $\Spec$ satisfies Lifted Monotonicity.  Complete each pending
exposure at the cut.  Every output at each resulting commitment-free stable
state is exposable by the consequence of Complete CALM above, giving
Availability.  The resulting participant projections are runs, so Complete
CALM gives their confluence and hence Partition Consistency.

Conversely, let any finite rows $(i_r,o_r)$ have a common refinement $j$.
Place their participants in separate partition components.  If there is only
one row, add a second participant initialized at $j$ and make no preparation
there.  Availability lets each original participant prepare and expose $o_r$
from $\gamma_0(i_r)$ without a commitment.  Partition Consistency then supplies
a
row $(j,u)$ surviving in every remainder table, so
$o_r\OutputOrd u$ for every $r$.  Form~(3) of
Proposition~\ref{prop:lifted-monotonicity-decomposition} gives Lifted
Monotonicity.
\end{proof}

Partition Consistency requires outputs from different components to be jointly
refinable, not identical or totally ordered; this is why the characterization
is bidirectional.  When Lifted Monotonicity fails, some permitted exposure
sequence either requires a commitment that cannot be established without a
cross-component wait, sacrificing Availability, or completes independently at
the cost of Partition Consistency.

A linearizable register is the classical instance of the tradeoff side: two
reads of concurrent writes may be explained by opposite serial orders, no
total order refines both, and the theorem withholds availability under
partition.  Whether a commitment requires agreement across a partition also
depends on which participant may hold which input and make which choice;
Appendix~\ref{app:deployment} refines the theorem by such a \emph{deployment}
and shows that a causal register, though it fails Lifted Monotonicity without
deployment restrictions, needs no coordination
when each participant answers only its own reads.  With the central results
in place, we now compare their consequences across the surrounding
literature.
